\documentclass{article}
\usepackage{authblk}
\usepackage{graphicx}
\usepackage{caption}
\usepackage{subcaption}

\title{A Comparison of the Rate of Reaction of $\rm{CaCO}_3$ Calcium Carbonate with Different concentrations of Hydrochloric Acid solutions}
\author[1]{Eduan Dervishi}
\affil[1]{\textit{Albanian College Tirana, Tirana, Albania}}
\date{March 23rd, 2026}
\begin{document}
\maketitle
\begin{abstract}
We present an investigation on how varying concentrations of $\rm{CaCO}_3$ reacted with hydrochloric acid ($\rm HCl$) influence the overall rate of reaction. By measuring the volume of $\rm {CO}_2$ gas produced, our findings suggest that an increase in reactant concentration correlates to an acceleration of initial reaction rate. The empirical data presented in this study validate the principles of collision theory, where a higher mass per given volume results in more successful molecular collisions, thereby resulting in more reactions per unit time. 
\end{abstract}
\section{Introduction}
Since time immemorial, the reaction between calcium carbonate and an acid has had vast medicinal and industrial applications, like in antacids, or home construction. The observation of carbon dioxide bubbles during the phenomenon of effervescence is also used to determine the presence of various carbonates in analytical chemistry as well as marine science. While the effects are well known, we have nevertheless set out to quantitatively investigate this relationship. 

\section{Methodology}
\subsection{Procedure}
Much like the premise and the nature of this investigation, its procedure is fairly simple. The steps taken are as follows:

\begin{enumerate}
\item{Measure two separate units of 0.50 g of $\rm{CaCO}_3$ powder using spatula.}

\item{Measure two units of $\rm 20.0\, cm^3$ at $\rm 1.00\, {dm}^3$ HCl solution using a 50 mL graduated cylinder, at 1.00 $\rm mol /{dm}^3$ and 0.3 $\rm mol /{dm}^3$ respectively.}
\item{Pour solutions into separate labelled flasks.}
\item{Submerge specified amount of $\rm{CaCO}^3$ in first HCl solution.}
\item{Repeat for other solution.}
\item{Using pressure sensor, record pressure in 5 second intervals for 30 seconds.}
\end{enumerate}

\subsection{Experimental Setup}
The experimental procedure in the section above was followed using this setup, where the separate HCl solutions have been excluded for redundancy.

\begin{figure}[!t]
\includegraphics[scale=0.3]{/Users/eduan/Downloads/setup111.png}
\centering
\textbf{\caption{Setup featuring solute, solvent, and relevant apparatus.}}
\end{figure}
\newpage

\section{Findings}
\subsection{Raw Data}
\begin{table} [h]
\begin{subtable}[c]{0.5\textwidth}
\centering
\begin{tabular}{|c|c|}
\hline
Time (s) & Pressure (kPa) \\
\hline
   0 & 112.26 \\

        5 & 114.83 \\

        10 & 116.25 \\
      
        15 & 117.21 \\
       
        20 & 117.86 \\

        25 & 117.99 \\
 
        30 & 116.96 \\
\hline
\end{tabular}
\subcaption{1.00 $mol / {dm}^3$}
\end{subtable}
\begin{subtable}[c]{0.5\textwidth}
\centering
\begin{tabular}{|c|c|}
\hline
 Time (s) & Pressure (kPa) \\
\hline
 0 & 104.09 \\
 5 & 107.11 \\
 10 & 108.65 \\
 15 & 109.65 \\
 20 & 110.78 \\
 25 & 111.68 \\
 30 & 112.45 \\
 \hline
 
\end{tabular}
\subcaption{0.30 $mol / {dm}^3$}
\end{subtable}
\end{table}

\begin{figure}[htbp!]
\includegraphics[scale=0.6]{/Users/eduan/Downloads/gaspressure.png}
\end{figure}
\section{Conclusion}


The data seems to suggest that the relationship between concentration and reaction rate is directly proportional. Note that there is a sharp increase in the rate of reaction - shown as a steep slope in the beginning - followed by a gradual decay until the rate completely flatlines. This is explained by particle kinetic theory. As the solution is saturated with products, there are less particles available for successful collision. This continues until there are none, and the reaction stops. 

However, there is far too little information from both the raw data and the graph to infer such a conclusion. Should we want a more robust result, more testing is required. With that being said, we compensate for our narrow dataset by noting that our results are consistent with the literature. 



\end{document}
